\section{Feasible outputs and completion of the allocation bound}
\label{sec:allocation-finish}

The proportion vectors now satisfy the constraints needed for the two
removal rules described in Section~\ref{sec:proportions}.  We define the
outputs, choose the absolute constants that preserve every child's lower
bound, and prove their movement estimate.  A hysteresis rule will keep
changes between the two formulas within the same movement budget.

All notation in this section is local to a star with parent \(v\), unless
a sum over stars is displayed.  In particular, \(S=\sum_i B_i\),
\(d_i=n_i+b_i\), and \(T=f_o^{\beta_v}\) when a dominant child \(o\)
exists.  Inactive coordinates, including their caps and outputs, are zero.
We extend proportion vectors by zero to the fixed child set when comparing
two times.  A marked child persists throughout its epoch; changes of its
identity, and transitions between empty and nonempty active sets, occur
only at wholesale steps.

\subsection{The output rules and the choice of constants}

The local data and proportion updates of
Sections~\ref{sec:allocation} and~\ref{sec:proportions} do not use the
current parent allocation.  Perform all of those updates first.  Once
\(q_v\) has been obtained in the top-down traversal, put
\begin{equation}
 R_0=\left(\sum_i d_i-q_v\right)_+.
 \label{eq:required-holes}
\end{equation}
If there is no active child, put every child output equal to zero.
Otherwise there are two possible rules:
\begin{align}
 \text{I:}\qquad &q_i=d_i-R_0\alpha_i\quad\text{for every active }i;
 \label{eq:output-rule-one}\\
 \text{II:}\qquad &q_i=d_i-Dw_i\quad(i\ne o\text{ active}),
 &q_o=\min\left\{d_o,q_v-\sum_{i\ne o}q_i\right\}.
 \label{eq:output-rule-two}
\end{align}
Rule~II is defined only when a dominant child exists.  With no dominant
child we always use Rule~I.  With a dominant child, the allowed ranges are
\begin{equation}
 \text{I requires }S\ge D/2,
 \qquad
 \text{II requires }S\le 2D.
 \label{eq:regime-ranges}
\end{equation}
At initialization and after a wholesale refresh, choose either rule in
its allowed range, with a fixed tie rule.  On other steps retain the old
rule unless it leaves its allowed range.  Thus an I-to-II switch occurs
only at \(S<D/2\), and an II-to-I switch only at \(S>2D\).  The ranges
cover all \(S,D\ge0\), including \(S=D=0\).

We specify the dependencies among the absolute constants before verifying
these formulas.  The rank functions, coarse size tolerances, and parameter
gaps give fixed constants
\[
 3\le\beta_v\le\beta_{\max},\qquad
 0\le h_i/\ell\le H_0,\qquad
 c_DT\le D\le C_DT,
\]
where the last comparison includes the assertion that \(D\) and \(T\)
vanish together.  Write \(c_g\) for the constant in
\eqref{eq:parameter-gap}.  At noncore ranks the parameter derivative is
bounded below by a fixed positive multiple of the rank function.
Consequently there is an absolute \(c_s>0\) such that, for a dominant
child,
\begin{equation}
 L_o^{\beta_v}-L_o^{\beta_o}
   \ge c_s\frac{u}{\ell}T.
 \label{eq:dominant-output-slack}
\end{equation}
Indeed the core contribution to this difference is
\((\beta_o-\beta_v)B_o\ge0\).  Affinity in the parameter and the
derivative comparison give
\(f_o^{\beta_v}-f_o^{\beta_o}
\ge c_f(\beta_o-\beta_v)f_o^{\beta_v}\), and the gap is at least
\(c_g u/\ell\).  This proof also covers \(T=0\) and \(u=0\).

Let \(C_z\) be an absolute constant for which
\eqref{eq:side-slack} reads
\(\sum_{i\ne o}\theta_i z_i\le C_z c_wu/\ell\).
Choose \(c_t,c_w,\epsilon_b>0\) so small that
\begin{gather}
 c_tH_0\le1,\qquad \beta_{\max}c_t\le c_g/4,
 \qquad c_wH_0\le1,\qquad c_w\le c_g/4,
 \label{eq:output-constant-smallness}\\
 C_zc_wC_D\le c_s/2,
 \qquad 4\epsilon_b/c_D\le c_g/4.
 \nonumber
\end{gather}
These choices involve only previously fixed absolute constants.  Next
choose the cap margin \(\lambda_b\), the parameter error \(\epsilon\),
and the fine tracking tolerance \(\rho\) as in Lemma~\ref{lem:caps},
decreasing \(\epsilon\) further so that \(\epsilon\le c_g/4\).
It follows that every side satisfies
\begin{equation}
 \beta^++\theta_i
 \le\beta_v+\frac{\epsilon+c_wh_i}{\ell}
 \le\beta_v+\frac{c_gh_i}{2\ell}
 \le\beta_i.
 \label{eq:side-parameter-room}
\end{equation}
Here \(h_i\ge1\) on regular coordinates.  The cap proof permits arbitrarily
small final choices of \(\epsilon,\rho\) after \(\lambda_b\) is fixed.
The bounded-domain constant \(C_w\) and the potential multipliers in
Section~\ref{sec:proportions} may now be taken sufficiently large.
For instance the required bound on \(C_w\) uses only
\(\beta^++\theta_i\le\beta_{\max}+1\).
This order fixes every local constant independently of \(k\), the degree,
the depth, and the separation ratio \(\tau\).  We choose \(\tau\) only
after proving the local movement estimate.

\begin{lemma}[Feasibility of the outputs]
\label{lem:output-feasibility}
Suppose \(q_v\ge L_v^{\beta_v}\).  The selected output rule satisfies
\[
 L_i^{\beta_i}\le q_i\le d_i\le C M_i\le C m_i,
 \qquad \sum_iq_i\le q_v,
\]
with an absolute constant \(C\).  In particular every child output is
nonnegative.  The conclusions include empty active sets, empty side sets,
and vanishing values of \(S,D,T\).
\end{lemma}

\begin{proof}
Superadditivity at the single parameter \(\beta_v\) gives
\begin{equation}
 q_v\ge L_v^{\beta_v}\ge\sum_iL_i^{\beta_v}.
 \label{eq:parent-lower-for-output}
\end{equation}
Inactive children have no core ranks and no positive lower-bound terms,
so their required lower bounds and their outputs are both zero.
For active children, expand
\(L_i^{\beta_v}=n_i-\beta_vB_i+f_i^{\beta_v}\).
Set \(e_o=b_o-T\) when a dominant child exists, and \(e_o=0\) otherwise.
Lemma~\ref{lem:caps} implies
\begin{equation}
 0\le e_o\le\epsilon_bT/\ell,
 \qquad R_0\le\beta_vS+e_o.
 \label{eq:required-hole-bound}
\end{equation}
To obtain the second inequality, subtract
\eqref{eq:parent-lower-for-output} from \(\sum_i d_i\): regular flex
caps contribute nonpositive errors, and only the dominant cap can
contribute the error \(e_o\).

First consider Rule~I with \(S>0\).  If there is a dominant child,
\eqref{eq:regime-ranges} and \(D\ge c_DT\) give
\begin{equation}
 T\le 2S/c_D,
 \qquad e_o/S\le 2\epsilon_b/(c_D\ell).
 \label{eq:rule-one-flex-comparison}
\end{equation}
The same bound on the error holds with error zero when there is no
dominant child.  At a regular child, \eqref{eq:alpha-feasibility} and
\eqref{eq:required-hole-bound} yield
\[
 R_0\alpha_i
 \le(\beta_v+e_o/S)(1+\eta_i)B_i.
\]
By \(\eta_i=c_th_i/\ell\), \(\eta_i\le1\), and
\eqref{eq:output-constant-smallness},
\begin{align*}
 (\beta_v+e_o/S)(1+\eta_i)-\beta_v
 &\le \beta_{\max}c_t h_i/\ell
       +4\epsilon_b/(c_D\ell)\\
 &\le c_gh_i/(2\ell)
 \le\beta_i-\beta_v.
\end{align*}
Thus \(R_0\alpha_i\le\beta_iB_i\), and the lower cap bound gives
\(q_i\ge n_i-\beta_iB_i+f_i^{\beta_i}=L_i^{\beta_i}\).
For the dominant coordinate use \(\alpha_o\le1\) on the error term:
\[
 R_0\alpha_o
 \le\beta_v(1+\eta_o)B_o+e_o
 \le\beta_oB_o+e_o.
\]
Here \(h_o=u\), so the same small choice of \(c_t\) works even for
arbitrarily small positive \(u\), and gives equality in the relevant
error when \(u=0\).  Since \(b_o-e_o=T\ge f_o^{\beta_o}\), this proves
the dominant lower bound as well.

If \(S=0\) in Rule~I, the presence of a dominant child forces \(D=0\)
by its allowed range, and therefore \(T=e_o=0\).  Without a dominant
child the error is zero by definition.  In either case
\eqref{eq:required-hole-bound} gives \(R_0=0\), so no division by \(S\)
is used: the outputs are their caps.  At regular children the cap lower
bound suffices, and at the dominant child \(B_o=T=0\) gives the claim.
For any nonempty active set, the simplex identity gives
\[
 \sum_iq_i=\sum_i d_i-R_0
          =\min\left\{\sum_i d_i,q_v\right\}\le q_v.
\]
Every hole is nonnegative, so \(q_i\le d_i\); the proved lower bound
ensures nonnegativity.

Now consider Rule~II.  Its side holes satisfy, by
\eqref{eq:side-feasibility} and \eqref{eq:side-parameter-room},
\[
 0\le Dw_i\le(\beta^++\theta_i)B_i\le\beta_iB_i.
\]
The side lower bounds follow from \(b_i\ge f_i^{\beta_i}\), exactly
by subtracting these holes from \(n_i+b_i\).  Write
\(q_o^{\mathrm{rem}}=q_v-\sum_{i\ne o}(d_i-Dw_i)\) for the
uncapped remaining amount.  The upper flex caps on the sides and
\eqref{eq:parent-lower-for-output} imply
\begin{align*}
 q_o^{\mathrm{rem}}
 &\ge L_o^{\beta_v}
        +\sum_{i\ne o}(-\beta_vB_i+Dw_i)\\
 &\ge L_o^{\beta_v}-D\sum_{i\ne o}\theta_i z_i\\
 &\ge L_o^{\beta_v}-C_zc_wC_D\frac{uT}{\ell}
 \ge L_o^{\beta_o}.
\end{align*}
The second line uses the lower inequality of
\eqref{eq:side-feasibility}, which applies because Rule~II has
\(S\le2D\), together with \(\beta^+\ge\beta_v\).  The last line
uses \eqref{eq:dominant-output-slack} and
\eqref{eq:output-constant-smallness}.  Independently,
\(d_o=n_o+b_o\ge n_o+f_o^{\beta_o}\ge L_o^{\beta_o}\).
Taking the smaller of these two admissible dominant amounts preserves
its lower bound.  It also gives the parent budget and the upper cap.

For completeness, when \(D=0\) in Rule~II its range forces \(S=0\),
and the comparison with \(T\) forces \(T=0\).  All side holes vanish
and \eqref{eq:required-hole-bound} gives \(R_0=0\); thus the dominant
cap also fits within the remaining budget.  When there are no sides,
\(u=0\), the side sums are empty, and both rules simply give
\(q_o=\min\{d_o,q_v\}\).  No side denominator or proportion is needed
in this case.  An empty active set was already assigned zero outputs.
Finally the cap estimate \(d_i\le CM_i\), together with
\(M_i\le m_i\), proves the claimed domination in every case.
\end{proof}

\subsection{Paying for changes of rule}

We next bound the cost of the hysteresis rule without assuming that the
absolute variation of the posterior sum \(S\) is small.  Fix a finite
horizon and let \(\mathcal W_v\) be the wholesale steps for this star.
The coarse charge \(C_{v,t}\) in \eqref{eq:coarse-data-budget} pays
core-flag changes, the variation of \(D\), and, on wholesale steps,
the old plus new active size totals.  All of its coefficients are
absolute.  Define the switch charge
\begin{equation}
 J_{v,t}=
 \begin{cases}
 D_t,&t\notin\mathcal W_v\text{ is an I-to-II switch},\\
 S_t,&t\notin\mathcal W_v\text{ is an II-to-I switch},\\
 0,&\text{otherwise}.
 \end{cases}
 \label{eq:switch-charge}
\end{equation}
There is a fixed finite constant \(D_v^{\mathrm{sw}}\), independent
of the hidden-process law and horizon, for which
\begin{equation}
 \E\sum_tJ_{v,t}\le C E_v+D_v^{\mathrm{sw}}.
 \label{eq:switch-budget}
\end{equation}

Here is the full potential accounting for this claim.  At the end of a
step use
\begin{equation}
 \Psi_v=
 \begin{cases}
 D,&\text{a dominant child exists and the rule is I},\\
 S,&\text{a dominant child exists and the rule is II},\\
 0,&\text{no dominant child exists}.
 \end{cases}
 \label{eq:switch-potential}
\end{equation}
Start every step by holding the old rule, mark, flags, and \(D\) fixed.
On this frozen state, the filtering change of \(\Psi_v\) is either
zero or the filtering change of the sum of the old flagged rank
posteriors.  Write \(F_{v,t}\) for this signed change.  The tower
property gives
\begin{equation}
 \E[F_{v,t}\mid\mathcal F_{t-1}]=0.
 \label{eq:switch-filtering-zero}
\end{equation}
On the subsequent drift leg, still with those old data fixed, the
increase is at most the sum of the absolute child-rank drifts at that
step.  Denote this nonnegative sum by \(e_{v,t}\); it has
\(\E\sum_t e_{v,t}\le E_v\).

On a non-wholesale step, changing the core flags changes the old-rule
potential by at most the number of changed flags, since each current
rank posterior lies in \([0,1]\).  Updating \(D\) changes it by at most
\(|\Delta D|\).  The other local updates leave this potential unchanged.
With the new inputs now in place, an I-to-II switch changes the potential
from \(D_t\) to \(S_t<D_t/2\), a drop of at least \(D_t/2\).
An II-to-I switch changes it from \(S_t\) to \(D_t<S_t/2\), a drop
of at least \(S_t/2\).

On a wholesale step we do not claim switch credit.  The old frozen
value of \(S\), even after filtering and drift, is at most the old
number of core flags, and that number is bounded by a constant times
the old active size total.  The old held \(D\) has the same size bound.
The newly chosen potential is likewise bounded by a constant times
the new active size total.  Thus the entire reset of this potential
costs at most \(C C_{v,t}\).  We obtain on every step the single
inequality
\begin{equation}
 \Psi_v(t)-\Psi_v(t-1)
 \le F_{v,t}+e_{v,t}+C C_{v,t}-\tfrac12 J_{v,t}.
 \label{eq:switch-one-step}
\end{equation}
Importantly, \(F_{v,t}\) is included on \emph{every} step before the
wholesale and switch tests; it is never multiplied by the indicator
of an event chosen from the new posteriors.  Taking expectations,
summing, and using \(\Psi_v(t)\ge0\) gives
\eqref{eq:switch-budget} from \eqref{eq:coarse-data-budget}.
The initial value is at most \(Ck\), uniformly over valid states.
This proves the asserted independence of law and horizon for this
additional constant.

\subsection{The local movement estimate}

\begin{lemma}[Movement on a star]
\label{lem:star-movement}
For the outputs above and a fixed star \(v\), there is a finite
\(D_v^{\mathrm{out}}\), depending only on its fixed finite structure
and \(k\), such that
\begin{equation}
 \E\sum_t\sum_{i\text{ child of }v}|q_i(t)-q_i(t-1)|
 \le C_{\mathrm{par}}\E\sum_t|q_v(t)-q_v(t-1)|
       +C\ell E_v+D_v^{\mathrm{out}}.
 \label{eq:star-output-movement}
\end{equation}
Here \(C_{\mathrm{par}}\) and \(C\) are absolute, independent of
\(k\), degree, depth, and \(\tau\).  In fact the output estimates
below allow \(C_{\mathrm{par}}=1\).
\end{lemma}

\begin{proof}
Compare the vectors of holes
\(\zeta_i=d_i-q_i\), adding cap variation at the end.  By
Lemma~\ref{lem:output-feasibility} these vectors are nonnegative.
In Rule~I one has
\begin{equation}
 \zeta=R_0\alpha,
 \qquad \|\zeta\|_1=R_0.
 \label{eq:rule-one-hole-identity}
\end{equation}
In Rule~II let \(\omega_i=Dw_i\) on the sides and
\(W_{\mathrm{side}}=\sum_{i\ne o}\omega_i\).  The definition of the
dominant output gives the exact identities
\begin{align}
 \zeta_i&=\omega_i\quad(i\ne o),&
 \zeta_o&=(R_0-W_{\mathrm{side}})_+,&
 \|\zeta\|_1&=\max\{R_0,W_{\mathrm{side}}\}.
 \label{eq:rule-two-hole-identity}
\end{align}
Indeed before replacing the signed excess by its positive part the
dominant hole equals
\((\sum_i d_i-q_v-W_{\mathrm{side}})_+\), and
\(W_{\mathrm{side}}\ge0\) allows that replacement.  These identities
remain valid when \(D\), \(S\), or the side set vanishes.

Suppress the time index and use superscripts \(-,+\) for old and new
step-end values.  On a non-wholesale step the dominant mark, if any,
is the same at both times.  First suppose both times use Rule~I and
the active sets are nonempty.  Adding and subtracting
\(R_0^+\alpha^-\), and using \(\|\alpha^-\|_1=1\), gives
\begin{equation}
 \|\Delta\zeta\|_1
 \le |\Delta R_0|+R_0^+\|\Delta\alpha\|_1
 \le |\Delta R_0|+C S^+\|\Delta\alpha\|_1.
 \label{eq:hole-movement-one}
\end{equation}
The last bound follows from \eqref{eq:required-hole-bound} and
\eqref{eq:rule-one-flex-comparison}, with the zero case proved in
Lemma~\ref{lem:output-feasibility}.  If the active sets are empty at
both times the hole change is zero.  A change between empty and
nonempty is wholesale and is treated below.

If both times use Rule~II, the positive-part map is 1-Lipschitz, so
\eqref{eq:rule-two-hole-identity} gives
\begin{align}
 \|\Delta\zeta\|_1
 &\le |\Delta R_0|+2\|\Delta\omega\|_1\\
 &\le |\Delta R_0|+2D^+\|\Delta w\|_1
                      +2C_w|\Delta D|.
 \label{eq:hole-movement-two}
\end{align}
For the first line, the change in the side-hole total is at most
\(\|\Delta\omega\|_1\); the other copy is the change of the side
vector itself.  For the second, use
\(\Delta(Dw)=D^+\Delta w+(\Delta D)w^-\) and
\(\|w^-\|_1\le C_w\).  Inserting or deleting a side causes no
extra term because these vectors have been extended by zero.

At an I-to-II switch, \(S^+<D^+/2\), so
\eqref{eq:required-hole-bound} and \(T^+\le D^+/c_D\) show
\(R_0^+\le C D^+\).  The new hole total is at most
\(\max\{R_0^+,C_wD^+\}\le C D^+\).  The old total is \(R_0^-\),
which is at most \(R_0^++|\Delta R_0|\).  Nonnegativity of the
two hole vectors therefore yields
\begin{equation}
 \|\Delta\zeta\|_1\le C D^++|\Delta R_0|.
 \label{eq:hole-downcross}
\end{equation}
At an II-to-I switch, \(S^+>2D^+\), whence
\(D^+\le S^+/2\) and \(R_0^+\le C S^+\).  The old total is
at most \(R_0^-+C_wD^-\), while the new total is \(R_0^+\).
Using \(R_0^-\le R_0^++|\Delta R_0|\) and
\(D^-\le D^++|\Delta D|\) gives
\begin{equation}
 \|\Delta\zeta\|_1
 \le C S^++C_w|\Delta D|+|\Delta R_0|.
 \label{eq:hole-upcross}
\end{equation}
These two arguments bound the full change across a switch without
requiring the unused rule to have held at the preceding inputs.

Finally, positive part is 1-Lipschitz in
\eqref{eq:required-holes}, and therefore
\begin{equation}
 |\Delta R_0|\le\sum_i|\Delta d_i|+|\Delta q_v|,
 \qquad
 \sum_i|\Delta q_i|\le\sum_i|\Delta d_i|+\|\Delta\zeta\|_1.
 \label{eq:cap-parent-variation}
\end{equation}
Combining all the non-wholesale cases gives one useful pointwise
bound, with a single copy of the parent movement:
\begin{align}
 \sum_i|\Delta q_i|
 \le{}& |\Delta q_v|+2\sum_i|\Delta d_i|\nonumber\\
 &+C\left(S^+\|\Delta\alpha\|_1
        +D^+\|\Delta w\|_1+|\Delta D|+J_{v,t}\right).
 \label{eq:output-pointwise-movement}
\end{align}
Only defined tracker terms are included.  In particular they are
zero for an empty active set, and the side term is zero without
sides or a dominant child.

At a wholesale step use the actual output domination directly:
\begin{equation}
 \sum_i|\Delta q_i|
 \le\sum_i(q_i^-+q_i^+)
 \le C\left(\sum_i a_i^-+\sum_i a_i^+\right)
 \le C C_{v,t}.
 \label{eq:wholesale-output-movement}
\end{equation}
There is no per-child or per-epoch initialization charge here; the
sum of sizes is exactly the wholesale quantity already paid.
Take expectations and sum \eqref{eq:output-pointwise-movement} over
non-wholesale steps and \eqref{eq:wholesale-output-movement} over
wholesale steps.  The fine cap budget \eqref{eq:fine-cap-budget},
the proportion budget \eqref{eq:proportion-movement}, the coarse
budget \eqref{eq:coarse-data-budget}, and \eqref{eq:switch-budget}
give \eqref{eq:star-output-movement}.  Cap movement is added once;
it is never used as an input to a potential having an \(O(\ell)\)
input slope.  The coefficient of the parent variation in
\eqref{eq:output-pointwise-movement} is one, which proves the
claimed absolute bound on \(C_{\mathrm{par}}\).
\end{proof}

\subsection{Uniform initialization and the tree sum}

We record why the finite constants just used are uniform over laws
and horizons.  For a fixed star there are finitely many rank trackers,
namely \(k\) ranks per child and the ranks used at the parent.
Each elementary reference potential is
\(p(1-p)+(p-p^\circ)^2\) with \(p,p^\circ\in[0,1]\), and hence
has a fixed bounded range.  Its tracking tolerance is fixed once
\(k\) is fixed.  This bounds all initial and terminal reference
allowances uniformly over the posteriors.  Superadditivity gives
\(\sum_i M_i\le k\); active size totals and core counts are at most
\(Ck\), and \(D\le Ck\).  The coarse refresh accounting thus has
a uniform finite initialization allowance as well.

For the proportion potentials, \(\|\alpha\|_1=1\) and
\(\|w\|_1\le C_w\).  The regular component estimates in
Section~\ref{sec:proportions} sum to \(C\ell k\), and the dominant
component estimate has that same bound uniformly as \(u\) tends
to zero.  The side-potential components sum to
\(C\ell(D\|w\|_1+S)\le C\ell k\).  Their signed terminal
values therefore have a uniform finite lower bound, as required
for their telescopes.  The switch potential is nonnegative and
bounded by \(Ck\).  These bounds use finitely many fixed tracker
coordinates and one pair of endpoints for each telescope, never
the number of histories, requests, or epochs.  Wholesale
reinitializations are paid jumps within those same telescopes.
Trackers remain attached to the fixed vertices even while their
coordinates are inactive; reusing a vertex label never creates a
new unpaid endpoint allowance.

Initialize the star data at time zero with the stated tracking
rules, run the proportion minimizations, and choose a valid regime.
Set \(q_{\rm root}=k\) and propagate the outputs top down.  The root
has \(L_{\rm root}^{\beta_{\rm root}}=k\), since its hidden count
is always \(k\).  Lemma~\ref{lem:output-feasibility} therefore
proves by induction, at initialization and after each step,
\[
 q_v\ge L_v^{\beta_v},\qquad
 0\le q_v\le C_0m_v,\qquad
 \sum_{i\text{ child of }v}q_i\le q_v
\]
for one absolute \(C_0\), enlarging it to cover the root.  At the
required leaf \(x\), its hidden count is at least one almost surely
under the current posterior.  Hence \(p_{x1}=0\),
\(l^{\beta_x}(0)=1\), and \(q_x\ge1\).  All choices have used
only the observed history and its conditional laws.

It remains to sum the movement estimate.  Multiply
\eqref{eq:star-output-movement} for an internal vertex \(v\) by
\(r_{|v|}\), and sum over the fixed finite tree.  Every nonroot
vertex appears exactly once as a child, with weight
\(r_{|v|-1}\), so the left-hand side is the quantity \(V\) in
Theorem~\ref{thm:allocation}.  The root has no movement.  At every
other internal vertex the parent term has weight
\(r_{|v|}=r_{|v|-1}/\tau\).  Thus
\begin{equation}
 V\le \frac{C_{\mathrm{par}}}{\tau}V
        +C\ell\sum_{v\text{ internal}}r_{|v|}E_v
        +D_{\mathrm{tree}},
 \qquad
 D_{\mathrm{tree}}=
    \sum_{v\text{ internal}}r_{|v|}D_v^{\mathrm{out}}<\infty.
 \label{eq:allocation-absorption}
\end{equation}
The rank-drift coupling \eqref{eq:rank-drift-tree} bounds the
middle sum by \(6P\), uniformly for \(\tau\ge2\).  Choose any
absolute \(\tau\ge\max\{2,2C_{\mathrm{par}}\}\); it may be
enlarged later to meet the independent rounding requirement.
Absorbing the first term yields
\[
 V\le 12C\ell P+2D_{\mathrm{tree}}.
\]
The constants multiplying movement have no dependence on the
depth, degree, number of leaves, or \(k\).  The finite additive
constant may depend on the fixed tree, its radii, and \(k\), but
the uniform endpoint bounds above show that it is independent
of the law and horizon.  This proves Theorem~\ref{thm:allocation}.
